\documentclass[11pt,a4paper]{article}

\usepackage[T1]{fontenc}
\usepackage[utf8]{inputenc}
\usepackage{lmodern}
\usepackage{amsmath,amssymb,mathtools}
\usepackage{geometry}
\usepackage{microtype}
\usepackage{enumitem}

\geometry{margin=2.5cm}
\setlength{\parindent}{0pt}
\setlength{\parskip}{0.6em}
\setlength{\emergencystretch}{2em}
\numberwithin{equation}{section}

\title{Continuum Scaling of the Link Weight\\
in a Four-Dimensional Lattice Path Sum}
\author{Sayed Heyder Ali Zia}
\date{}

\begin{document}

\maketitle

\section{Question and scope}

We consider an infinite four-dimensional Euclidean
hypercubic lattice with spacing $\Delta$. We assign to
each path $\gamma$ from the origin to a lattice point
$x$ the weight
\begin{equation}
    w(\gamma)=e^{-m_0L(\gamma)},
\end{equation}
where we use $L(\gamma)$ for the number of links in
the path.

We seek the dependence $m_0=m_0(\Delta)$ that keeps
long paths relevant as $\Delta\to0$. We focus on this
dependence rather than on evaluating the full propagator.

\section{Assumptions and notation}

We make the following assumptions:
\begin{enumerate}[label=(\roman*)]
    \item We work on an infinite Euclidean lattice.
    \item We allow only nearest-neighbour steps.
    \item We allow backtracking, repeated sites,
    repeated links, and self-intersections.
    \item We count different ordered sequences of
    steps as different paths.
    \item We assign equal weights to all directions.
    \item We sum over all finite path lengths,
    including the zero-link path at the origin.
    \item We keep the physical endpoint separation
    fixed as we take the continuum limit.
\end{enumerate}

We write the lattice sites as
\begin{equation}
    x=\Delta n,
    \qquad
    n=(n_1,n_2,n_3,n_4)\in\mathbb Z^4.
\end{equation}
We distinguish the physical position $x$ from the
dimensionless integer coordinate $n$.

More precisely, we take a sequence of lattice points
$x_\Delta=\Delta n_\Delta$ approaching a fixed continuum
point $x$. We suppress the additional subscript when
it causes no ambiguity.

We denote the number of links by $N$ and the physical
length by $\ell$:
\begin{equation}
    N=L(\gamma),
    \qquad
    \ell=N\Delta.
\end{equation}
Since $N$ is dimensionless, we identify $m_0$ as a
dimensionless cost per link. We reserve $m$ for the
continuum mass, which has dimensions of inverse length.

We use units $\hbar=c=1$. We adopt the prescribed
Euclidean weight as the physical input:
\begin{equation}
    A(\gamma)=m_0N,
    \qquad
    w(\gamma)=e^{-A(\gamma)}.
\end{equation}

We take $m_0$ to be real and seek a stable continuum
limit with $m^2\geq0$.

\section{Path counting and the residual length weight}

We denote the four coordinate unit vectors by
$e_\mu$, with $\mu=1,2,3,4$. In lattice units, we
choose each step from
\begin{equation}
    \{\pm e_1,\pm e_2,\pm e_3,\pm e_4\}.
\end{equation}
We therefore have eight possible directions at every
step.

If we do not restrict the endpoint, we count $8^N$
paths with $N$ links. We obtain the total weight
\begin{equation}
    8^Ne^{-m_0N}
    =
    e^{-(m_0-\ln8)N}.
\end{equation}

We define
\begin{equation}
    \delta(\Delta)=m_0(\Delta)-\ln8,
    \qquad
    q(\Delta)=8e^{-m_0(\Delta)}=e^{-\delta(\Delta)}.
    \label{eq:delta}
\end{equation}
We interpret $\delta$ as the residual cost per step
after accounting for the exponential number of paths.

We now let $C_N(n)$ denote the number of $N$-step
paths from the origin to the specified endpoint $n$.
We write
\begin{equation}
    \mathcal G_\Delta(\Delta n;m_0)
    =
    \sum_{N=0}^{\infty}C_N(n)e^{-m_0N}.
\end{equation}

We introduce
\begin{equation}
    P_N(n)=\frac{C_N(n)}{8^N}.
\end{equation}
We use $P_N(n)$ as a counting tool: we recognize it
as the endpoint probability for an unbiased walk
whose eight step directions each have probability
$1/8$.

We obtain the exact rewriting
\begin{equation}
    \boxed{
    \mathcal G_\Delta(\Delta n;m_0)
    =
    \sum_{N=0}^{\infty}q^NP_N(n).
    }
    \label{eq:walksum}
\end{equation}

We have separated the endpoint constraint, contained
in $P_N(n)$, from the additional suppression of long
paths, contained in $q^N$.

\section{Critical value and convergence}

We identify the critical value
\begin{equation}
    \boxed{m_{0,c}=\ln8.}
\end{equation}

For fixed lattice endpoint $n$, we use the large-$N$
estimate
\begin{equation}
    P_N(n)\sim cN^{-2},
    \qquad c>0,
\end{equation}
along the allowed parity subsequence.

We understand the power $N^{-2}$ from diffusion:
we spread over distances of order $\sqrt N$ in
four directions and therefore over a volume of
order $(\sqrt N)^4=N^2$.

We must respect the parity restriction
\begin{equation}
    N\equiv\sum_{\mu=1}^{4}n_\mu\pmod2.
\end{equation}
We obtain no paths for the opposite parity because
each step changes the parity of the sum of the
coordinates.

We distinguish three cases:
\begin{enumerate}
    \item For $m_0<\ln8$, we have $q>1$. We obtain
    exponential growth that overwhelms the
    power-law decay of $P_N(n)$, so the sum diverges.

    \item For $m_0>\ln8$, we have $q<1$. We obtain
    exponential suppression of long paths and a
    convergent sum.

    \item For $m_0=\ln8$, we have $q=1$. We still
    obtain a finite fixed-endpoint sum in four
    dimensions because $\sum_NN^{-2}$ converges.
\end{enumerate}

We therefore retain $m_0=\ln8$ as an exception to
the statement that every fixed value leads either
to divergence or to the suppression of long-path
continuum behaviour.

We also distinguish convergence at one endpoint
from convergence after summing over every endpoint.
At criticality, we have
\begin{equation}
    \sum_n\mathcal G_\Delta(\Delta n;\ln8)
    =
    \sum_{N=0}^{\infty}\sum_nP_N(n)
    =
    \sum_{N=0}^{\infty}1
    =
    \infty.
\end{equation}

\section{The meaning of long paths}

We first consider a shortest path to $n$. We need
\begin{equation}
    N_{\min}=\sum_{\mu=1}^{4}|n_\mu|
\end{equation}
links.

At fixed nonzero physical separation, we have
$N_{\min}\sim\Delta^{-1}$. We therefore cannot
characterize the desired regime merely by saying
that $N$ becomes large: even a shortest path has
infinitely many links in the continuum limit.

We instead seek fluctuating paths on the diffusive
scale. We write their displacement after $N$
unbiased steps as
\begin{equation}
    X_N=\Delta\sum_{j=1}^{N}s_j.
\end{equation}
We have
\begin{equation}
    \langle s_j\rangle=0,
    \qquad
    |s_j|^2=1,
    \qquad
    \langle s_j\cdot s_k\rangle=0
    \quad(j\neq k).
\end{equation}
We consequently obtain
\begin{equation}
    \langle |X_N|^2\rangle=N\Delta^2.
\end{equation}

We identify the scale on which a fluctuating walk
explores a finite physical distance as
\begin{equation}
    \boxed{
    N\Delta^2=\text{finite},
    \qquad
    N\sim\Delta^{-2}.
    }
    \label{eq:diffusion}
\end{equation}
We use this relation as a scaling statement, not as
a claim that every contributing path has the same
length.

We also find
\begin{equation}
    \ell=N\Delta\sim\Delta^{-1}.
\end{equation}
We therefore allow divergent microscopic path
lengths while keeping the endpoints a finite
physical distance apart.

For fixed $m_0>\ln8$, we can make the loss of
finite-distance propagation explicit. Since
$P_N(n)\leq1$ and $q<1$, we have
\begin{equation}
    0\leq\mathcal G_\Delta(\Delta n;m_0)
    \leq
    \sum_{N=N_{\min}}^\infty q^N
    =
    \frac{q^{N_{\min}}}{1-q}.
\end{equation}
We obtain exponential decay as $\Delta\to0$ at
fixed nonzero separation.

We do not infer that exactly shortest paths alone
contribute. We regard the shortest-path wording
as a heuristic description, since longer paths
can have substantial multiplicity.

\section{Dependence of $m_0$ on $\Delta$}

We retain from \eqref{eq:walksum} the residual
length factor
\begin{equation}
    q^N=e^{-\delta(\Delta)N}.
    \label{eq:residual}
\end{equation}

We consider the diffusive limit
\begin{equation}
    N\Delta^2\longrightarrow u,
    \qquad 0<u<\infty.
    \label{eq:u}
\end{equation}

We rewrite the exponent as
\begin{equation}
    \delta(\Delta)N
    =
    \frac{\delta(\Delta)}{\Delta^2}
    \bigl(N\Delta^2\bigr).
    \label{eq:exponent}
\end{equation}

We obtain a definite finite limiting suppression
by requiring
\begin{equation}
    \frac{\delta(\Delta)}{\Delta^2}
    \longrightarrow\alpha,
    \qquad 0\leq\alpha<\infty.
    \label{eq:alpha}
\end{equation}

We then find
\begin{equation}
    q^N\longrightarrow e^{-\alpha u}.
    \label{eq:finiteweight}
\end{equation}

Since $\delta(\Delta)=m_0(\Delta)-\ln8$, we obtain
\begin{equation}
    \boxed{
    m_0(\Delta)
    =
    \ln8+\alpha\Delta^2+o(\Delta^2),
    \qquad \alpha\geq0.
    }
    \label{eq:main}
\end{equation}

\subsection{Explanation of Equation \eqref{eq:main}}

We obtain \eqref{eq:main} directly from the limit
in \eqref{eq:alpha}. We can express that limit by
writing
\[
    \frac{\delta(\Delta)}{\Delta^2}
    =
    \alpha+r(\Delta),
    \qquad
    r(\Delta)\longrightarrow0.
\]
We multiply by $\Delta^2$:
\[
    \delta(\Delta)
    =
    \alpha\Delta^2+\Delta^2r(\Delta).
\]
We denote $\Delta^2r(\Delta)$ by $o(\Delta^2)$
because its ratio to $\Delta^2$ vanishes. We then
substitute $\delta=m_0-\ln8$ and obtain
\eqref{eq:main}.

We can also understand the result physically.
We have $N\sim\Delta^{-2}$ for the long fluctuating
paths of interest. We therefore require the
residual cost per link to be of order $\Delta^2$,
so that its product with the number of links
remains finite.

We assign $\alpha$ dimensions of inverse length
squared. We consequently keep
$\alpha\Delta^2$ dimensionless, consistently with
the dimensions of $m_0$.

We additionally require
$m_0(\Delta)\geq\ln8$ at sufficiently small finite
spacing. We cannot approach the critical value
from below while retaining the unrestricted
positive path sum, because the sum would diverge
at each such spacing.

\subsection{Identification of the continuum mass}

We identify the coefficient $\alpha$ using the
one-step momentum dependence.

We let $p$ denote physical Euclidean momentum and
define
\begin{equation}
    \varphi_\Delta(p)
    =
    \left\langle e^{ip\cdot\Delta s}\right\rangle
    =
    \frac18\sum_{\mu=1}^{4}
    \left(e^{i\Delta p_\mu}
    +e^{-i\Delta p_\mu}\right).
\end{equation}
We use
$p\cdot x=\sum_{\mu=1}^{4}p_\mu x_\mu$ and
$p^2=\sum_{\mu=1}^{4}p_\mu^2$.

We expand at fixed $p$:
\begin{equation}
    \varphi_\Delta(p)
    =
    \frac14\sum_{\mu=1}^{4}\cos(\Delta p_\mu)
    =
    1-\frac{\Delta^2p^2}{8}
    +O(\Delta^4|p|^4).
\end{equation}

We combine this factor with the residual
length weight:
\begin{equation}
    \log\left(q^N\varphi_\Delta(p)^N\right)
    =
    -N\delta(\Delta)
    +N\log\varphi_\Delta(p).
\end{equation}

Using $N\Delta^2\to u$ and
$\delta(\Delta)=\alpha\Delta^2+o(\Delta^2)$, we obtain
\begin{equation}
    q^N\varphi_\Delta(p)^N
    \longrightarrow
    \exp\left[-u\left(\alpha+\frac{p^2}{8}\right)\right].
\end{equation}
We have discarded terms proportional to
$N\Delta^4$, which vanish because
$N\Delta^4=(N\Delta^2)\Delta^2\to0$.

We define the proper-time parameter
\begin{equation}
    s=\frac{u}{8}
    =
    \lim_{\Delta\to0}\frac{N\Delta^2}{8}.
\end{equation}
We then write the limiting factor as
\begin{equation}
    e^{-s(p^2+8\alpha)}.
\end{equation}

We identify the continuum mass through
\begin{equation}
    m^2=8\alpha.
\end{equation}
We therefore obtain
\begin{equation}
    \boxed{
    m_0(\Delta)
    =
    \ln8+\frac{m^2\Delta^2}{8}+o(\Delta^2).
    }
    \label{eq:mass}
\end{equation}

We identify $s$ as a parameter with dimensions
of length squared. We distinguish it from both
the link count $N$ and the physical path length
$N\Delta$.

We may choose
\begin{equation}
    m_0(\Delta)=\ln(8+m^2\Delta^2)
\end{equation}
as one function satisfying \eqref{eq:mass}.
We do not require this particular finite-spacing
choice; we require its asymptotic behaviour.

\subsection{Limiting cases}

We distinguish the following possibilities:
\begin{enumerate}
    \item If
    $\delta(\Delta)/\Delta^2\to+\infty$,
    we suppress paths with any fixed positive
    value of $N\Delta^2$. We lose the desired
    finite-mass diffusive limit.

    \item If
    $\delta(\Delta)/\Delta^2\to\alpha>0$,
    we retain long paths with finite suppression.
    We obtain a finite positive mass squared
    $m^2=8\alpha$.

    \item If
    $\delta(\Delta)/\Delta^2\to0$ with
    $\delta(\Delta)\geq0$,
    we retain the diffusive regime without a
    residual mass suppression. We obtain the
    massless scaling limit.

    \item If $\delta(\Delta)<0$, we encounter
    a divergent unrestricted path sum at that
    spacing.
\end{enumerate}

We therefore require both an approach to the
critical value and control over the rate of that
approach. We cannot determine the continuum
behaviour from $m_0(\Delta)\to\ln8$ alone.

\section{Final result}

We find the general finite-suppression scaling
\begin{equation}
    \boxed{
    m_0(\Delta)-\ln8
    =
    \alpha\Delta^2+o(\Delta^2),
    \qquad \alpha\geq0,
    }
\end{equation}
with $m_0(\Delta)\geq\ln8$ at sufficiently small
finite spacing.

For a specified continuum mass $m$, we write
\begin{equation}
    \boxed{
    m_0(\Delta)
    =
    \ln8+\frac{m^2\Delta^2}{8}+o(\Delta^2).
    }
\end{equation}

We interpret $\ln8$ as the compensation for the
exponential number of available walks. We
interpret the correction of order $\Delta^2$
as the residual suppression that remains finite
for paths with $N\sim\Delta^{-2}$.

We retain $m_0=\ln8$ as the critical massless
case. We also retain the distinction between
the survival of long-path contributions and
the separate normalization of a continuum
propagator.

\subsection*{AI use declaration}

\textbf{Model(s) used:}
We used ChatGPT (Codex; the exact underlying model
version was not available in the conversation).\\

\textbf{Revisions/corrections made to model output,
and clarifications added:}
We made many revisions and clarifications,
including the distinction between link count
and physical length, the dimensions of $m_0$,
the fixed-separation continuum limit, the
critical exception $m_0=\ln8$, the interpretation
of long paths, the qualification of
shortest-path domination, and the derivation
of Equation~\eqref{eq:main}.\\

\textbf{Comparisons/checks against the lecture notes:}
We made two comparisons with the relativistic
path-integral discussion in Chapter~6,
printed pages~52--55:
(1) the cancellation between path multiplicity
and the exponential length weight;
(2) the diffusive scaling in which the number
of steps times the squared step length remains
finite. We accounted for the different
conventions for link count and physical length.\\

\textbf{Calculation steps, limiting cases, or
consistency checks verified by hand:}
We report one independent hand check.

\textbf{Calculation steps, limiting cases, or
consistency checks subjected to automated
verification (via Mathematica, Python, LLMs):}
We used one LLM-assisted review of the scaling
argument and its algebraic consistency.
We did not perform an independent Mathematica
or Python verification.

\end{document}